\documentclass[11pt]{article}
\usepackage[utf8]{inputenc}
\usepackage[T1]{fontenc}
\usepackage{amsmath}
\usepackage{amssymb}
\usepackage{geometry}
\usepackage[style=numeric,backend=biber]{biblatex}
\addbibresource{bib/references.bib}
\geometry{margin=1in}

\title{Variable Definitions from Derrida and Gardner (1984)}
\date{}

\begin{document}

\maketitle

\section*{Variables and Definitions}

\begin{itemize}
    \item $\lambda$: The parameter used for the weak disorder expansion ($\lambda \rightarrow 0$).
    \item $E$: A fixed energy parameter.
    \item $\gamma(E)$ or $\gamma$: The Lyapounov exponent associated with a given product of random matrices.
    \item $\psi_n$: The wave function evaluated at site $n$.
    \item $V_n$: The random potential at site $n$.
    \item $\rho(V)$: The given probability distribution of the random potential $V$.
    \item $U_n$: A two-component vector defined as containing the wave functions $\psi_{n+1}$ and $\psi_n$.
    \item $N$: The upper limit in the product of random matrices or the total number of terms.
    \item $\epsilon$: A parameter that is $0$ for integrable dynamical systems and non-zero for mixing properties. It also represents an infinitesimal imaginary part added to real values of energy.
    \item $R_n$: The ratio of wave functions at adjacent sites, defined as $\psi_n / \psi_{n-1}$, which measures the direction of the vector $U_n$.
    \item $q$: A real parameter used to define an energy $E \neq 2 \cos(qa)$ that does not belong to the spectrum of the pure system.
    \item $A, B_n, C_n, D_n$: Variables in the weak disorder expansion of $R_n$ that do not depend on $\lambda$.
    \item $x$: A scaled parameter defining the deviation from an energy level, such as $E-2 = \lambda^{4/3}x$ near the band edge or $E' = 2 \cos(\pi\alpha) + \lambda^2 x$.
    \item $P(R, E, \lambda)$: The stationary probability distribution for the variable $R_n$.
    \item $\tilde{\rho}(E)$: The density of states at energy $E$.
    \item $t$: A scaled variable for $R$, defined by the relationship $R = 1 + \lambda^{2/3}t$.
    \item $H(t, x, \lambda)$: A scaled probability distribution function derived from $P$ via the transformation $H(t, x, \lambda) = \lambda^{2/3}P(1+\lambda^{2/3}t, 2+\lambda^{4/3}x, \lambda)$.
    \item $H_0, H_1, H_2$: Successive functions in the power expansion of $H(t, x, \lambda)$ with respect to $\lambda$.
    \item $C, C_1$: Arbitrary constants of integration.
    \item $X$: A scaled parameter defined as $x / \langle V^2 \rangle^{2/3}$.
    \item $\alpha$: A rational parameter used to define specific energy points via $E = 2 \cos(\pi\alpha)$.
    \item $\varphi$: An angular variable introduced through the change of variables $R = \sin(\varphi+\pi\alpha)/\sin\varphi$.
    \item $G(\varphi)$: A transformed probability distribution defined as $P(R, E, \lambda) \frac{dR}{d\varphi}$.
    \item $G_0, G_1, G_2$: Successive functions in the $\lambda$ expansion of $G(\varphi)$.
    \item $E'$ : A perturbed energy value in the neighborhood of $2 \cos(\pi\alpha)$ defined as $E' = E + \lambda^2 x$.
    \item $r, s$: Integers used to define the rational number $\alpha = r/s$.
    \item $W(\varphi)$: A periodic function of period $\pi/3$ that appears when solving the equations in the neighborhood of $E=1$.
\end{itemize}

\printbibliography

\end{document}